\section{The ordinary-dilate model has a different limit}\label{sec:ordinary}

For clarity about the basis in \cite{nbdraft}, define
\[
 d^{\mathrm{ord}}_M=\inf_c\Big\|1-\sum_{k\le M}c_k\{kx\}\Big\|_{L^2(0,1)}.
\]
This is the ordinary integer-dilate problem. The exact formula below is independent of the location of zeta zeros.

\begin{proposition}
Let \(J_2(d)=d^2\prod_{p\mid d}(1-p^{-2})\) and \(\varphi\) denote Euler's totient. Then
\[
 (d_M^{\mathrm{ord}})^2=\frac{1}{1+3S_M},
 \quad S_M=\sum_{d\le M}\frac{\varphi(d)^2}{J_2(d)},\qquad
 S_M\sim M\prod_p\left(1-\frac{2}{p(p+1)}\right).
\]
\end{proposition}
\begin{proof}
The centered fractional part has nonzero Fourier coefficients \(-1/(2\pi i \ell)\) at \(\ell\in\mathbb Z\setminus\{0\}\). Parseval therefore gives
\[
 \int_0^1\{jx\}\{kx\}\dd x=\frac14+\frac{\gcd(j,k)^2}{12jk}.
\]
Indeed, after writing \(j=dj_0,k=dk_0\) with \((j_0,k_0)=1\), the matching nonzero frequencies are \(j k_0 r=k j_0 r\). The corresponding coefficient products sum to \((j_0k_0)^{-1}\sum_{r\ne0}(4\pi^2r^2)^{-1}=1/(12j_0k_0)\). The constant Fourier coefficient contributes \(1/4\).

Put \(u_d=\sum_{d\mid k\le M}c_k/k\) and \(t=\sum_{k\le M}c_k\). The map from the \(c_k\) to the \(u_d\) is invertible by descending triangular elimination. The divisor identities \(\sum_{d\mid k}\varphi(d)=k\) and \(\sum_{d\mid k}J_2(d)=k^2\) follow by multiplicativity and a telescoping sum at each prime power. Thus \(t=\sum_d\varphi(d)u_d\), and the squared residual is
\[
 |1-t/2|^2+\frac{1}{12}\sum_{d\le M}J_2(d)|u_d|^2.
\]
For fixed \(t\), weighted Cauchy--Schwarz minimizes the last sum to \(|t|^2/S_M\) before the factor \(1/12\), attained by \(u_d=t\varphi(d)/(S_MJ_2(d))\). The remaining quadratic in \(t\) is minimized at \(t=6S_M/(1+3S_M)\) and has the displayed value.

Finally \(a(d)=\varphi(d)^2/J_2(d)\) is multiplicative with \(a(p^k)=(p-1)/(p+1)\) for \(k\ge1\). Thus \(a=1*b\), where \(b\) is supported on squarefree integers, \(b(p)=-2/(p+1)\), and \(b(1)=1\). Since \(\sum|b(d)|/d=\prod_p(1+2/[p(p+1)])<\infty\), dividing \(\sum_{d\le M}b(d)\lfloor M/d\rfloor\) by \(M\) and using dominated convergence gives the stated product. Every factor is positive and their deviations from \(1\) are summable.
\end{proof}

The actual Mellin transform of \(\{kx\}\) on \((0,1)\) is
\[
 \int_0^1 \{kx\}x^{s-1}\dd x
 =\frac{k}{s+1}-\frac{k-1}{s}+\frac{k^{-s}}s\sum_{r=1}^{k-1}r^s
 \quad(\R s>0).
\]
This follows by writing \(\lfloor kx\rfloor=\sum_{r=1}^{k-1}\one_{[r/k,1]}(x)\) almost everywhere and integrating. The transform cannot be replaced by a multiple of \(k^{-s}\zeta(s)/s\). For example, at \(k=1\) it is \(1/(s+1)\). Thus the \(M^{-1/2}\) distance scale proved above is about the ordinary dilates. The balanced reciprocal functions have the different Mellin image of Lemma~\ref{lem:nbmellin}.
